%==============================================================================
% Capítulo 2: Marco Teórico
%==============================================================================
\section{Marco teórico}\label{sec:marco-teorico}

Este capítulo establece los fundamentos conceptuales sobre los cuales se construye la investigación. Se abordan ocho ejes temáticos que proporcionan el andamiaje teórico necesario para comprender el diseño, la implementación y la evaluación del sistema propuesto.

\subsection{Gestión de servicios de tecnologías de la información}\label{subsec:gestion-servicios-ti}

La gestión de servicios de tecnologías de la información (ITSM) constituye una disciplina sistematizada que articula procesos, personas y tecnologías con el propósito de entregar valor al negocio a través de servicios tecnológicos confiables y mensurables. \textcite{Galup2009_itsm} caracterizan la disciplina como un cuerpo de prácticas orientado a alinear la operación tecnológica con los objetivos del negocio, donde la mesa de ayuda opera como punto único de contacto entre el usuario y el área de tecnología. Este enfoque rompe con la concepción anterior que entendía a la informática como un área de soporte estrictamente reactivo y la posiciona, en cambio, como un proveedor de servicios sujeto a acuerdos de nivel de servicio explícitos.

Dentro de este marco general, la gestión de incidentes constituye uno de los procesos operativos centrales. Su objetivo principal consiste en restaurar el servicio afectado en el menor tiempo posible, minimizando el impacto adverso en la operación del negocio. El proceso comprende cuatro etapas secuenciales: la recepción y registro inicial, la clasificación y priorización, la asignación al área responsable y, finalmente, la resolución y cierre. \textcite{Galup2009_itsm} advierten que los procesos manuales de soporte tienden a degradarse cuando el volumen de solicitudes crece de manera no lineal o cuando el personal debe alternar entre tareas operativas heterogeneas, situación que coincide con la observada en el entorno empresarial analizado en el presente trabajo.

\subsection{Mesa de ayuda y modelo ITIL}\label{subsec:mesa-ayuda-itil}

El marco de buenas prácticas ITIL (Information Technology Infrastructure Library), en su edición de operaciones de servicio \parencite{OGC2011_itil}, publicada por The Stationery Office (TSO) por encargo del Cabinet Office y cuyo titular histórico es el Office of Government Commerce, establece la distinción conceptual entre incidente ---definido como una interrupción no planificada del servicio o una reducción de su calidad--- y solicitud de servicio ---definida como un pedido formal de información, asesoramiento, acceso a un componente o ejecución de una operación estandarizada---. Esta distinción resulta especialmente relevante para el diseño del clasificador propuesto en este trabajo, ya que orienta la categorización inicial y condiciona la priorización subsiguiente.

El marco ITIL postula adicionalmente tres principios operativos que guian el diseño de la solución propuesta. El primero es la centralización de la recepción mediante un punto único de contacto (SPOC, por su sigla en inglés), que en el sistema propuesto se materializa a través del flujo unificado de orquestación que normaliza las entradas de los tres canales. El segundo es la operación bajo acuerdos de nivel de servicio (SLA), que cuantifican la calidad esperada en términos de tiempo de respuesta y tiempo de resolución. El tercero es la mejora continua del proceso, que en el sistema propuesto se habilita mediante la trazabilidad de cada decisión del clasificador y la retroalimentación de los sectores responsables.

Los modelos tradicionales de mesa de ayuda se basan en la interacción manual entre el usuario y el operador, lo que implica tiempos de respuesta variables y una dependencia directa del personal disponible. Cuando el volumen de incidentes excede la capacidad del equipo o cuando los operadores deben alternar simultáneamente entre tareas operativas y de soporte, la calidad del servicio se degrada de manera observable. La automatización del registro inicial constituye, en consecuencia, una intervención que aborda el cuello de botella en el punto de mayor sensibilidad del proceso, sin requerir la reestructuración completa del modelo operativo.

\subsection{Automatización de procesos y orquestación de flujos}\label{subsec:automatizacion-orquestacion}

Se entiende aquí por automatización de procesos empresariales (BPA, por su sigla en inglés) la reducción de la intervención humana en tareas repetitivas mediante la ejecución automática de acciones definidas a partir de eventos específicos. \textcite{Russell2021_ai_modern} ubican a la automatización inteligente dentro del paradigma de los agentes orientados a objetivos, en el cual un sistema reactivo percibe eventos del entorno, decide acciones según una política preestablecida y actúa modificando el estado del entorno. Esta caracterización ofrece un marco conceptual robusto para el diseño de sistemas que combinan reglas determinísticas con decisiones derivadas de modelos probabilísticos, como el clasificador híbrido propuesto en este trabajo.

Se entiende aquí por plataformas de integración como servicio (iPaaS) las herramientas de orquestación que posibilitan integrar distintos servicios, procesar información y ejecutar acciones automáticas en función de condiciones predefinidas. Estas plataformas se caracterizan por ofrecer modelos visuales de construcción de flujos, soporte nativo para webhooks y disparadores reactivos, y capacidad de transformación de datos entre sistemas heterogeneos. \textcite{n8n2024_docs} documenta las capacidades de N8N, la herramienta seleccionada para este trabajo, que se distingue dentro del segmento por su modelo de licencia de código fuente disponible y su capacidad de despliegue completamente autoalojado, característica que resulta decisiva en escenarios donde el tratamiento de datos sensibles exige mantener la información dentro del perimetro organizacional.

El paradigma de orquestación se diferencia conceptualmente del paradigma de coreografía. \textcite{Hohpe2003_eip} establecen que la orquestación centraliza la lógica de control en un componente único, lo que simplifica la trazabilidad de cada ejecución pero introduce un punto de coordinación cuya disponibilidad condiciona al sistema completo. La coreografía, en cambio, distribuye la coordinación entre los participantes mediante eventos publicados en un bus comun, lo que aumenta la robustez frente a fallos parciales pero dificulta la observabilidad. El presente trabajo adopta el enfoque de orquestación por simplicidad operativa y por la prioridad otorgada a la trazabilidad auditable de cada decisión del sistema, en línea con los requisitos de la Ley~25.326 \parencite{Ley25326_2000}.

\subsection{Human-in-the-loop y grados de autonomía}\label{subsec:human-in-the-loop}

El paradigma \emph{human-in-the-loop} (HITL) designa los esquemas en los que una persona interviene en el ciclo de decisión de un sistema automático para etiquetar, corregir, validar o supervisar sus salidas. \textcite{MosqueiraRey2023_human_in_the_loop} sistematizan el aprendizaje con intervención humana y muestran que las anotaciones y correcciones incorporadas al ciclo mejoran la calidad del modelo cuando los datos escasean o el dominio es ambiguo; \textcite{Zanzotto2019_human_in_the_loop} sostiene que el diseño de estos sistemas debe reservar a la persona un rol activo en el lazo ante la ambigüedad semántica y la responsabilidad sobre las consecuencias de la decisión. \textcite{Amershi2019_human_ai} formulan guías de interacción que exigen transparencia sobre la confianza del sistema y control del usuario, \textcite{Shneiderman2020_human_centered_ai} propone combinar alta automatización con supervisión significativa, y \textcite{Wu2022_ai_chains} muestran que la descomposición de tareas con puntos de control humano hace al sistema inspeccionable.

Los grados de autonomía se ordenan según la proporción de decisiones delegadas y la granularidad de la revisión. La \emph{revisión total} somete cada salida a validación humana; la \emph{revisión selectiva} delega las decisiones de alta confianza y deriva a la persona solo las que caen por debajo de un umbral; y la \emph{operación autónoma con auditoría} ejecuta sin intervención previa y conserva un registro trazable para revisión posterior. El sistema propuesto adopta el grado intermedio: clasifica de forma automática y escala a revisión humana todo caso cuya confianza declarada cae por debajo del umbral, según la regla \texttt{requiere\_revision\_humana = (confianza < 0{,}70)}, evaluada en la \autoref{sec:resultados}. La calibración de esa confianza condiciona la cobertura efectiva de la salvaguarda (\autoref{sec:discusion}).

\subsection{Procesamiento de lenguaje natural y modelos de lenguaje grandes}\label{subsec:nlp-llm}

El procesamiento de lenguaje natural (NLP) constituye un área de la inteligencia artificial orientada a la interpretación del lenguaje humano por parte de sistemas informáticos. \textcite{Jurafsky2023_speech_language} sistematizan la disciplina en torno a tres ejes principales: el modelado estadístico de secuencias, la representación distribuida de palabras y la comprensión semántica mediante arquitecturas neuronales profundas. Mediante estas técnicas resulta posible analizar texto o audio transcrito, identificar palabras clave, determinar la intención del usuario y clasificar contenido en categorías predefinidas.

La evolución del campo durante la última decada ha estado dominada por la arquitectura Transformer \parencite{Vaswani2017_attention}, cuyo mecanismo de atención permite a los modelos capturar dependencias de largo alcance sin las limitaciones de las redes neuronales recurrentes. El mecanismo de atención multi-cabeza, en particular, permite que el modelo pondere simultáneamente distintas partes de la secuencia de entrada, lo que resulta especialmente útil para la clasificación de descripciones textuales donde la información relevante puede aparecer en cualquier posición del texto.

Sobre esta arquitectura se construyeron los modelos de lenguaje grandes (LLM) contemporáneos, los cuales exhiben capacidades robustas de clasificación en escenarios de pocos ejemplos (\emph{few-shot}) o sin ejemplos (\emph{zero-shot}), fenómeno conocido como aprendizaje en contexto. \textcite{Brown2020_few_shot} demostraron que modelos suficientemente grandes pueden generalizar a tareas nuevas a partir de instrucciones expresadas en lenguaje natural, sin requerir ajuste fino sobre datos específicos del dominio. Esta capacidad es particularmente relevante para el contexto de esta investigación, donde el volumen de datos etiquetados resulta insuficiente para entrenar un clasificador supervisado tradicional desde cero.

La forma en que se fórmula la instrucción condiciona el resultado tanto como la capacidad del modelo. El aprendizaje en contexto admite dos configuraciones principales: el planteamiento \emph{zero-shot}, en el que la instrucción describe la tarea sin ejemplos, y el \emph{few-shot}, en el que se incorporan unos pocos casos resueltos dentro del propio mensaje para orientar la respuesta. La evaluación multilingüe de \textcite{ahuja2023mega} y el ajuste por instrucciones de \textcite{lai2023okapi} evidencian que la selección de ejemplos, su orden y la redacción de la instrucción modifican el rendimiento de manera sustancial, por lo que el diseño del \emph{prompt} constituye una decisión de ingeniería y no un detalle accesorio.

Los LLM generan, junto con la etiqueta asignada, una estimación de confianza que no necesariamente coincide con la probabilidad empírica de acierto. La calibración designa el grado de correspondencia entre la confianza declarada por el modelo y su corrección observada: un modelo bien calibrado acierta con una frecuencia próxima a la confianza que expresa. \textcite{xuan2025mmluprox} documentan, a partir de una evaluación comparativa en múltiples idiomas, que la calibración depende del modelo y del dominio. En el sistema propuesto este concepto fundamenta el umbral operativo de revisión humana, ya que las clasificaciones con confianza inferior a 0,70 se derivan a supervisión; la fiabilidad de ese umbral depende tanto de la exactitud como de la calibración del clasificador.

Una limitación central de los modelos generativos es la alucinación, es decir, la producción de contenido plausible en su forma pero falso o no sustentado en la entrada \parencite{Ji2023_hallucination}. En una tarea de clasificación la alucinación puede manifestarse como la asignación de una categoría inexistente o como una justificación inventada a partir de indicios débiles. \textcite{colavito2025benchmarking} muestran que el etiquetado automático mediante LLM requiere validación contra un conjunto de referencia y no puede aceptarse sin control. Esta observación motiva dos decisiones del sistema propuesto: la restricción del espacio de salida a las cinco categorías canónicas y la validación de cada clasificación contra el conjunto de etiquetas válidas antes de su persistencia.

El rendimiento de los LLM tampoco es uniforme entre idiomas ni entre variedades de un mismo idioma. \textcite{zhang2023multilingual} documentan que los modelos multilingües no generalizan de manera homogenea fuera del inglés, y la evaluación de variedades del español peninsular y latinoamericano de \textcite{grandury2025laleaderboard} y \textcite{mayorrocher2026spanish} muestra diferencias de desempeño atribuibles al dialecto y a la disponibilidad de recursos. Esta sensibilidad resulta directamente relevante para el presente trabajo, cuyo corpus está redactado en español rioplatense, una variedad escasamente representada en los conjuntos de evaluación internacionales.

La evaluación de un LLM no se agota en una medida única de exactitud. Los benchmarks multilingües como los de \textcite{xuan2025mmluprox} y \textcite{ahuja2023mega} combinan múltiples tareas y reportan resultados desagregados, en reconocimiento de que una cifra global puede ocultar un desempeño dispar entre categorías. Esta advertencia, que este capítulo desarrolla en la Sección~\ref{sec:metricas-multiclase} para el caso de clasificadores multiclase, fundamenta la decisión metodológica de reportar, además de la exactitud global, métricas por clase y el valor F1 macro promediado.

En el contexto de una mesa de ayuda, el uso de NLP permite clasificar incidentes y derivarlos al sector correspondiente sin intervención manual, mejorando la velocidad de respuesta y reduciendo la carga operativa del personal encargado del registro. La elección entre un clasificador supervisado entrenado sobre corpus específico, un modelo de lenguaje grande de propósito general invocado mediante instrucciones, o un esquema híbrido que combine ambas estrategias, constituye una decisión de diseño que depende de la disponibilidad de datos etiquetados, del presupuesto de inferencia disponible y de los requerimientos de soberanía sobre los datos procesados.

\subsection{Reconocimiento automático del habla}\label{subsec:reconocimiento-habla}

Los sistemas de reconocimiento automático del habla (ASR) permiten convertir audio en texto mediante modelos entrenados sobre grandes corpus de habla anotada, cuyo marco clásico describen \textcite{Rabiner1993_speech_recognition}, mientras que la generación contemporánea se apoya en redes neuronales profundas cuyo exponente reciente con alta robustez multilingüe es el modelo Whisper
\parencite{Radford2023_whisper}. El sistema propuesto no integra un motor de ASR: la transcripción efectiva del canal telefónico la provee un servicio externo (\autoref{subsec:integracion-canal-telefonico}), de modo que la noción de reconocimiento automático del habla se conserva únicamente como marco conceptual de ese canal.

\subsection{Bases de datos relacionales y persistencia transaccional}\label{subsec:bd-relacionales}

La persistencia de los incidentes en una base de datos relacional permite garantizar las propiedades de atomicidad, consistencia, aislamiento y durabilidad (ACID) propias del modelo transaccional clásico \parencite{Date2003_database}. PostgreSQL, el motor seleccionado para este trabajo, ofrece soporte completo del estándar SQL, extensiones para tipos de datos avanzados como JSON binario y arreglos, y un modelo de concurrencia basado en control multiversional (MVCC) que resulta adecuado para cargas mixtas de lectura y escritura \parencite{PostgreSQL2024_docs}.

La elección de un motor relacional se comprende mejor al contrastarla con alternativas concretas. Un motor de documentos como MongoDB almacena registros con esquema flexible y favorece la evolución de la estructura sin migraciones, pero no impone las relaciones entre incidentes, sectores, estados y canales de origen, de modo que la integridad referencial queda delegada a la aplicación. Un motor clave-valor en memoria como Redis ofrece latencias de acceso muy bajas y resulta adecuado para cache o contadores efímeros, pero no está orientado a la persistencia duradera ni a consultas relacionales complejas. El dominio de este trabajo exige precisamente esas propiedades: las relaciones entre incidentes, sectores, estados y canales se modelan como tablas vinculadas mediante claves foráneas con integridad referencial declarativa, las operaciones de registro y actualización requieren las garantías ACID de la transacción clásica, el dashboard de analítica se apoya en consultas relacionales de agregación sobre múltiples tablas y la trazabilidad de cada decisión del sistema demanda un registro duradero y auditable. Esta combinación de requisitos justifica la elección de PostgreSQL, en línea con los requisitos normativos identificados en la sección de consideraciones éticas y legales (\autoref{sec:aspectos-legales}).

\subsection{Métricas de evaluación de clasificadores multiclase}
\label{sec:metricas-multiclase}

La evaluación de clasificadores multiclase requiere métricas que capturen tanto el desempeño global como el desempeño por clase. \textcite{Powers2011_evaluation} describe el conjunto canónico de métricas derivadas de la matriz de confusión, entre las cuales destacan la exactitud global (\emph{accuracy}), la precisión por clase, la sensibilidad o exhaustividad (\emph{recall}) por clase, y la medida F1 como media armónica de las dos anteriores.

La exactitud global se define como la proporción de predicciones correctas sobre el total de casos evaluados:
\[
\text{Exactitud} = \frac{1}{N} \sum_{i=1}^{N} \mathbb{1}[\hat{y}_{i} = y_{i}]
\]
donde \(N\) es el número de casos, \(\hat{y}_{i}\) la clase predicha para el caso \(i\), \(y_{i}\) su clase verdadera y \(\mathbb{1}[\cdot]\) la función indicadora, que vale uno cuando la predicción coincide con la etiqueta y cero en caso contrario. Esta definición es válida directamente para cualquier número de clases y es la que subyace a la matriz de confusión de cinco sectores de la \autoref{sec:resultados}. No se emplea la formulación binaria basada en verdaderos y falsos positivos y negativos, que presupone una única clase de interés y resulta inadecuada para un problema multiclase.

\textcite{Sokolova2009_performance_measures} advierten que en escenarios con clases desbalanceadas la exactitud global puede resultar enganosa y recomiendan reportar valores macro promediados, los cuales otorgan igual peso a cada clase con independencia de su frecuencia. El valor F1 macro promediado se define como:
\[
F1_{\text{macro}} = \frac{1}{k} \sum_{i=1}^{k} F1_{i}
\]
donde \(k\) es el número de clases y \(F1_{i}\) es el valor F1 de la clase \(i\). El presente trabajo adopta esta recomendación y reporta tanto exactitud global como F1 macro y métricas desagregadas por clase, permitiendo al lector evaluar la robustez del clasificador frente a la distribución específica del corpus.

Adicionalmente, la matriz de confusión constituye el instrumento estándar de presentación de resultados en clasificación multiclase, ya que permite identificar patrones específicos de error: confusiones sistemáticas entre dos categorías particulares, sesgos hacía una clase mayoritaria o aciertos preferenciales sobre una clase específica. La inspección de la matriz orienta el análisis cualitativo posterior y constituye la base para cualquier propuesta de mejora del clasificador.

\newpage
